\chapter{Óptica Ondulatoria}
\label{ch:ondulatoria}


La visión geométrica de la luz, si bien elegante y suficiente para muchas de las aplicaciones elementales, no capta por completo la verdadera naturaleza de la luz, pues en la naturaleza la luz no se mantiene confinada en un rayo mientras se propaga por el espacio libre, sino que tiende a ``esparcirse'' hacia las vecindades, entre otros fenómenos que la teoría geométrica no explica ni contempla.

En concreto, esta manera geométrica de entender la luz se vio enfrentada a serios problemas para explicar fenómenos como la interferencia, la difracción y la birrefringencia, por mencionar algunos. Con el tiempo se fueron acumulando experiencias que contradecían el modelo geométrico, hasta que las falencias de este modelo se hicieron innegables.

Por otro lado, aparentemente inconexo en un principio, las investigaciones concernientes a los fenómenos eléctricos y magnéticos dieron fruto a grandes avances en el entendimiento de estas fuerzas a distancia observadas en cuerpos ``terrenales'' como la magnetita. Fue Faraday quien entendió cómo se influían mutuamente, lo que sugería una misma naturaleza, y también observó experimentalmente cómo la luz era influida por estas fuerzas, ahora mejor entendidas. James Clerk Maxwell, en una tarea verdaderamente unificadora y excelsa matemáticamente, sintetizó las ideas de su época en las ecuaciones que llevan su nombre y, por una necesidad física de conservación y el conocimiento matemático de las ecuaciones que la aseguraban, fue más allá que sus predecesores postulando las ondas como consecuencia de respetar este principio de conservación para la carga eléctrica \cite{maxwell1865}.

Las ondas electromagnéticas de Maxwell tenían que viajar a cierta velocidad que resultó ser la misma que la velocidad a la que viajaban las ondas luminosas. 
Finalmente, se llegó a concluir que la luz visible era radiación electromagnética. 



\section{Ecuaciones de Maxwell}

En su forma general, las ecuaciones de Maxwell pueden escribirse como \cite{jackson1999}




\begin{align}
\nabla \cdot \mathbf{E}(\mathbf x, t) &= \frac{\rho(\mathbf x, t)}{\epsilon_0} 
& \qquad 
\nabla \cdot \mathbf{B}(\mathbf x, t) &= 0, \\
\nabla \times \mathbf{E}(\mathbf x, t) &= -\,\frac{\partial \mathbf B(\mathbf x, t)}{\partial t}
& \qquad
\nabla \times \mathbf B(\mathbf x, t) &= \mu_0\,\mathbf J(\mathbf x, t) 
+ \mu_0\epsilon_0\,\frac{\partial \mathbf E(\mathbf x, t)}{\partial t},
\label{eq:maxwell_equaitons}
\end{align}

y, en el caso particular del \emph{vacío}, donde no existen cargas ni corrientes libres,
\[
\rho(\mathbf x,t) = 0, \qquad \mathbf J(\mathbf x,t) = \mathbf 0,
\]
las ecuaciones toman la forma


\begin{align}
\nabla \cdot \mathbf{E}(\mathbf x, t) &= 0 
& \qquad 
\nabla \cdot \mathbf{B}(\mathbf x, t) &= 0, 
\\
\nabla \times \mathbf{E}(\mathbf x, t) &= -\,\frac{\partial \mathbf B(\mathbf x, t)}{\partial t}
& \qquad
\nabla \times \mathbf B(\mathbf x, t) &= \mu_0\epsilon_0\,\frac{\partial \mathbf E(\mathbf x, t)}{\partial t},
\label{eq:maxwell_equation_vacuum}
\end{align}


De la combinación adecuada de estas ecuaciones se concluye que los campos satisfacen ecuaciones de onda de la forma

\begin{align}
\nabla^2 \mathbf E(\mathbf x, t) - \frac{1}{c^2}\frac{\partial^2 \mathbf E(\mathbf x, t)}{\partial t^2} &= 0,
\label{eq:wave_equation_E} 
\\
\nabla^2 \mathbf B(\mathbf x, t) - \frac{1}{c^2}\frac{\partial^2 \mathbf B(\mathbf x, t)}{\partial t^2} &= 0,
\label{eq:wave_equation_B}
\end{align}
donde
\[
c = \frac{1}{\sqrt{\mu_0 \epsilon_0}}
\]
es la velocidad característica de propagación de la onda electromagnética en el vacío.


Por el teorema de Fourier sabemos que toda función admite una representación armónica, de modo que podemos escribir el campo eléctrico en dicha representación como sigue

\begin{equation}
    \vb{E}(\vb x, t) = \frac{1}{2\pi}\int_{-\infty}^{\infty} \vb{E}_\omega(\vb x) e^{-i\omega t} d\omega
    \label{eq:E_field_harmonic_representation}
\end{equation}

En la representación armónica el campo magnético se obtiene fácilmente, como 


\begin{equation}
    \vb{B}_\omega = \frac{1}{i\omega} \nabla \times \vb{E}_\omega,
    \label{eq:B_harmonic_func_of_E_harmonic}
\end{equation}

por lo que en lo sucesivo omitiremos el campo magnético en nuestro análisis, pues la información del campo eléctrico nos es suficiente. De este modo, cada componente armónica satisface

\begin{equation}
    \nabla^2 \vb{E}(\vb x) + k^2 \vb E = 0, 
    \label{eq:E_field_wave_equation}
\end{equation}

donde $k = \omega/c$, la ecuación fundamental para la propagación del campo armónico. 

La ecuación \ref{eq:E_field_wave_equation} es equivalente a una ecuación de Helmholtz escalar para cada una de sus componentes en una base independiente del espacio; sea $\Psi(\vb x)$ una de esas componentes, escribimos

\begin{equation}
    \nabla^2 \Psi (\vb x) + k^2 \Psi (\vb x) = 0.
    \label{eq:Helmholtz_equation}
\end{equation}


Que es la ecuación base de los \textit{métodos escalares de propagación}. Siendo los métodos de propagación escalares, en últimas, soluciones a la ecuación escalar de Helmholtz dadas unas condiciones de contorno. Los métodos se dividen en exactos y aproximados, según si
resuelven la ecuación de Helmholtz \eqref{eq:Helmholtz_equation} de forma exacta o aproximada.


% ======================
\section{Propagación de los campos}

\subsection{Método de espectro angular}

El método de espectro angular permite determinar el campo escalar en un plano $z>0$ a partir de su valor conocido en el plano inicial $z=0$. Su fundamento consiste en descomponer el campo como una superposición continua de ondas planas, cada una de las cuales evoluciona de manera simple bajo la ecuación de Helmholtz.

Para ello, se separa primero el operador laplaciano en su parte transversal y longitudinal,

\begin{equation}
\nabla^2
=
\nabla_\perp^2
+
\frac{\partial^2}{\partial z^2},
\qquad
\nabla_\perp^2
=
\frac{\partial^2}{\partial x^2}
+
\frac{\partial^2}{\partial y^2},
\label{eq:laplacian_transverse_longitudinal}
\end{equation}

de modo que la ecuación de Helmholtz toma la forma

\begin{equation}
\left(
\nabla_\perp^2
+
\frac{\partial^2}{\partial z^2}
+
k^2
\right)
\Psi(\vb r,z)
=
0,
\label{eq:helmholtz_split}
\end{equation}

donde $\vb r=(x,y)$ representa la coordenada transversal.

Introducimos ahora la transformada de Fourier bidimensional sobre el plano transversal mediante

\begin{equation}
\mathscr F_\perp\{\Psi\}
=
\tilde{\Psi}(\mathbf k_\perp;z)
=
\iint_{\mathbb R^2}
\Psi(\vb r,z)\,
e^{-i\,\mathbf k_\perp\cdot\vb r}
\,d^2\vb r,
\label{eq:fourier_forward}
\end{equation}

y su inversa

\begin{equation}
\Psi(\vb r,z)
=
\frac{1}{(2\pi)^2}
\iint_{\mathbb R^2}
\tilde{\Psi}(\mathbf k_\perp;z)\,
e^{i\,\mathbf k_\perp\cdot\vb r}
\,d^2\mathbf k_\perp,
\label{eq:fourier_inverse}
\end{equation}

donde $\mathbf k_\perp=(k_x,k_y)$ es el vector de onda transversal.

Aplicando $\mathscr F_\perp$ sobre la ecuación \eqref{eq:helmholtz_split}, y usando que

\begin{equation}
\mathscr F_\perp
\left\{
\nabla_\perp^2\Psi
\right\}
=
-|\mathbf k_\perp|^2
\tilde{\Psi},
\label{eq:fourier_laplacian}
\end{equation}

se obtiene la ecuación diferencial ordinaria

\begin{equation}
\frac{\partial^2\tilde{\Psi}}{\partial z^2}
+
\left(
k^2
-
|\mathbf k_\perp|^2
\right)
\tilde{\Psi}
=
0.
\label{eq:spectral_ode}
\end{equation}

Esta expresión muestra que cada componente espectral transversal evoluciona independientemente como una onda plana.

Definiendo la componente longitudinal del vector de onda mediante la relación de dispersión

\begin{equation}
k^2
=
|\mathbf k_\perp|^2
+
k_z^2,
\label{eq:dispersion_relation}
\end{equation}

se obtiene

\begin{equation}
k_z(\mathbf k_\perp)
=
\begin{cases}
\sqrt{
k^2-|\mathbf k_\perp|^2
},
&
|\mathbf k_\perp|^2\leq k^2,
\\[6pt]
i\sqrt{
|\mathbf k_\perp|^2-k^2
},
&
|\mathbf k_\perp|^2>k^2.
\end{cases}
\label{eq:kz_definition}
\end{equation}

La elección de la rama con $\operatorname{Im}(k_z)\geq0$ garantiza que las componentes evanescentes no crezcan exponencialmente para $z>0$.

Las componentes tales que
$|\mathbf k_\perp|^2\leq k^2$
corresponden a ondas propagantes, mientras que aquellas con
$|\mathbf k_\perp|^2>k^2$
representan campos evanescentes que decaen exponencialmente con la distancia.

La solución general de \eqref{eq:spectral_ode} es

\begin{equation}
\tilde{\Psi}(\mathbf k_\perp;z)
=
A(\mathbf k_\perp)\,
e^{ik_z z}
+
B(\mathbf k_\perp)\,
e^{-ik_z z},
\label{eq:general_solution_spectral}
\end{equation}

y para propagación hacia adelante ($z>0$) se conserva únicamente la componente físicamente saliente,

\begin{equation}
\tilde{\Psi}(\mathbf k_\perp;z)
=
\tilde{\Psi}(\mathbf k_\perp;0)\,
e^{ik_z(\mathbf k_\perp)\,z}.
\label{eq:spectral_propagation}
\end{equation}


Así el espectro transversal se propaga retrasando cada componente acorde a la fase $\left( \sqrt{k^2 - k_r^2} \right ) z$, nótese cómo cada componente se retrasa de manera distinta según $k_r$, los mayores retardos corresponden a los menores valores que puede tomar $k_r$, es decir las frecuencias transversales cercanas a cero.  






Finalmente, sustituyendo \eqref{eq:spectral_propagation} en la transformada inversa \eqref{eq:fourier_inverse}, obtenemos





\begin{equation}
\Psi(\vb r,z)
=
\frac{1}{(2\pi)^2}
\iint_{\mathbb R^2}
\tilde{\Psi}_0(\mathbf k_\perp)\,
\exp
\left[
i\left(
\mathbf k_\perp\cdot\vb r
+
k_z(\mathbf k_\perp)\,z
\right)
\right]
\,d^2\mathbf k_\perp.
\label{eq:angular_spectrum_propagation}
\end{equation}



En esta representación el campo $\Psi$ en el plano $z$ es una suma de ondas planas con vector de onda $\vb k = \vb k_\perp + k_z(\vb k_\perp)\,\uvec z$, cuya amplitud es el espectro angular del campo en $z=0$. 





De manera operacional, la propagación puede escribirse como

\begin{equation}
\Psi(\vb r,z)
=
\mathscr R_{ASM}(z)\,
\Psi(\vb r,0),
\label{eq:operator_form}
\end{equation}

donde

\begin{equation}
\mathscr R_{ASM}(z)
=
\mathscr F_\perp^{-1}
\left\{
e^{i\,k_z(\mathbf k_\perp)\,z}
\mathscr F_\perp \left\{ \right\} \right\}
\label{eq:propagation_operator}
\end{equation}


O equivalente, 

\begin{equation}
\mathscr R_{ASM}(z)
=
\mathscr F_\perp^{-1}
\left\{
  e^{i\, k \sqrt{1 - (k_r/k)^2 }  \,z}
\mathscr F_\perp \left\{ \right\} \right\}
\label{eq:propagation_operator_explicit}
\end{equation}





El operador $\mathscr R_{ASM}(z)$ diagonaliza el laplaciano transversal $\nabla_\perp^2$, transformando la ecuación diferencial parcial de Helmholtz en una operación espectral.








% ==== SEC : LÍMITE DE RESOLUCIÓN ===
\subsection{Límite de resolución} 
\label{sec:resolution_limit}

Tomamos como punto de partida el resultado del método de espectro angular, ecuación \eqref{eq:angular_spectrum_propagation}, que expresa el campo como superposición de ondas planas propagantes y evanescentes. En particular, el campo escalar en el medio puede escribirse como:

\begin{equation}
\begin{aligned}
\Psi(\mathbf r)
&= \frac{1}{(2\pi)^2}
\int_{-\infty}^{\infty}\!\!\int_{-\infty}^{\infty}
\tilde\Psi(k_x,k_y,0)\,
e^{i(k_x x + k_y y + k_z z)}\, dk_x\,dk_y .
\end{aligned}
\end{equation}

donde

\begin{equation}
k_z = \sqrt{k^2 - k_x^2 - k_y^2},
\qquad
k = \lvert\mathbf k\rvert = \frac{2\pi n}{\lambda},
\end{equation}

con \(n\) el índice de refracción del medio y \(\lambda\) la longitud de onda en el vacío.

La dirección de propagación de cada onda plana viene dada por

\begin{equation}
\frac{\mathbf k\cdot\hat{\mathbf z}}{k} = \cos\theta
\qquad\Rightarrow\qquad
\frac{k_z}{k} = \cos\theta .
\end{equation}

De donde se obtiene

\begin{equation}
\sin\theta
= \frac{k_\perp}{k}
= \frac{\sqrt{k_x^2 + k_y^2}}{k},
\end{equation}

con \(k_\perp = \sqrt{k_x^2 + k_y^2}\).

La apertura numérica del sistema está dada por \cite{bornwolf1999}

\begin{equation}
\mathrm{NA} = n \sin\theta.
\label{eq:def_NA}
\end{equation}


La finitud de los sistemas ópticos, debida a la finitud de las lentes y de las pupilas, impone
un ángulo máximo de aceptación desde el cual los rayos del objeto pueden ser emitidos para que pasen al sistema, 
sea $\alpha$ este ángulo, se sigue que las componentes que se emiten con un ángulo $\theta$ deben ser menores o iguales al ángulo máximo de aceptación $\alpha$, de lo que se sigue que


\begin{equation}
n\sin\alpha \ge n\sin\theta = \mathrm{NA}.
\end{equation}

\begin{figure}[H]
    \centering
    \begin{tikzpicture}[>=stealth,
    ray/.style={color=black!80!black, thick},
    optaxis/.style={dash dot, gray!80}]

    % --- Parámetros geométricos ---
    \def\distLente{5}
    \def\altLente{2}
    \def\altObjeto{1.5}

    % --- Coordenadas ---
    \coordinate (O) at (0,0);
    \coordinate (L) at (\distLente,0);
    \coordinate (TopL) at (\distLente, \altLente);
    \coordinate (EjeEnd) at (\distLente + 2, 0);

    % --- Eje óptico ---
    \draw[optaxis] (-1,0) -- (EjeEnd) node[right, black] {Eje óptico};

    % --- Objeto (plano) ---
    \draw[thick] (0,-\altObjeto) -- (0,\altObjeto);
    \fill[gray!10] (-0.15,-\altObjeto) rectangle (0,\altObjeto);
    \node[below=2pt] at (0,-\altObjeto) {Objeto};

    % --- Sistema óptico (lente) ---
    \begin{scope}[shift={(L)}]
        \draw[fill=gray!5, thick]
            (0, \altLente) .. controls (-0.4, 0.5) and (-0.4, -0.5) .. (0, -\altLente)
            .. controls (0.4, -0.5) and (0.4, 0.5) .. (0, \altLente) -- cycle;
        \node[below=2.2cm] at (0,0) {Sistema Óptico};
    \end{scope}

    % --- Rayo marginal y ángulo alpha ---
    \draw[ray, ->] (O) -- (TopL);

    \coordinate (Ref) at (1,0);
    \pic [draw, "$\alpha$", angle eccentricity=1.4, angle radius=1.2cm] {angle = Ref--O--TopL};

    \fill (O) circle (1.5pt);

\end{tikzpicture}

    \caption{Ilustración del ángulo de apertura $\alpha$.}
    \label{fig:na_geometry}
\end{figure}


Que es equivalente a


\begin{equation}
k_\perp \le k \sin\theta_{\max}
= k\,\frac{\mathrm{NA}}{n}.
\end{equation}

En términos de frecuencias espaciales

\begin{equation}
\begin{aligned}
f_\perp &= \frac{k_\perp}{2\pi}
\le \frac{k}{2\pi}\,\frac{\mathrm{NA}}{n}
= \frac{n/\lambda}{1}\,\frac{\mathrm{NA}}{n}
= \frac{\mathrm{NA}}{\lambda}.
\end{aligned}
\end{equation}

Por tanto, el sistema transmite frecuencias espaciales hasta

\begin{equation}
f_{\max} = \frac{\mathrm{NA}}{\lambda}.
\end{equation}


Considérese ahora un objeto limitado a una extensión \(d\) en la dirección \(x\). Su desarrollo en serie de Fourier contiene armónicas con números de onda

\begin{equation}
\begin{aligned}
k_x &= \frac{n\pi}{d},
\\
n &= 0,1,2,\dots
\end{aligned}
\end{equation}

y por tanto

\begin{equation}
f_x = \frac{k_x}{2\pi} = \frac{n}{2d},
\end{equation}

Para que la estructura correspondiente sea resoluble, al menos el primer armónico \(n=1\) debe ser transmitido por el sistema:

\begin{equation}
\frac{1}{2d} \le f_{\max} = \frac{\mathrm{NA}}{\lambda}.
\end{equation}

De aquí se obtiene el límite de resolución de Abbe:

\begin{equation}
d \ge \frac{\lambda}{2\,\mathrm{NA}}.
\end{equation}

Por lo que el límite inferior de detalle resoluble está dado por

\begin{equation}
  d_{min} = \frac{\lambda}{2\,\mathrm{NA}} ,
\label{eq:abbe_resolution_limit}
\end{equation}

donde \(\lambda\) es la longitud de onda en el vacío, y \(\mathrm{NA}\) es la apertura numérica definida en la ecuación \eqref{eq:def_NA}.

\begin{figure}[H]
    \centering
    \thesisincludegraphics[width=0.85\textwidth]
      {figures/Loss_of_detail_in_small_apertures.jpg}
      {figures/fast/Loss_of_detail_in_small_apertures.jpg}
    \caption{Pérdida de detalle al reducir la apertura numérica: las frecuencias espaciales altas quedan fuera de banda y no son transmitidas por el sistema óptico. Se observa que las frecuencias espaciales más altas contienen los detalles más finos.}
    \label{fig:loss_detail_small_apertures}
\end{figure}

De acá se sigue que aumentar el tamaño de la pupila y aumentar el índice de refracción $n$ da lugar a una mayor resolución, por lo que las aproximaciones paraxiales limitan la resolución que los sistemas pueden alcanzar debido a que no contemplan frecuencias mayores, donde están los detalles más finos.



\section{Solución por la Segunda Identidad de Green}


Además del tratamiento en el dominio de las frecuencias espaciales mediante el método del espectro angular, la propagación en el vacío puede formularse también desde una perspectiva puramente espacial. En lugar de descomponer el campo en ondas planas, esta formulación reconstruye el campo en un punto del espacio a partir de los valores que el campo toma sobre una superficie que encierra el dominio de interés.

La base matemática de esta formulación es la segunda identidad de Green \cite{arfken2013}. Su importancia radica en que permite transformar el problema diferencial asociado a la ecuación de Helmholtz en una representación integral de frontera. En una región homogénea, isotrópica y libre de fuentes, cada componente cartesiana del campo eléctrico satisface una ecuación escalar de Helmholtz. Por tanto, si \(U(\mathbf r)\) representa una componente escalar del campo, se tiene

\begin{equation}
\nabla^2 U(\mathbf r)+k^2U(\mathbf r)=0.
\label{eq:scalar_Helmholtz_equation}
\end{equation}

Sea \(V\) un volumen limitado por una superficie cerrada \(\Sigma(V)\), y sea \(G(\mathbf r,\mathbf r')\) la función de Green saliente asociada al operador de Helmholtz, definida por

\begin{equation}
\left(\nabla'^2+k^2\right)G(\mathbf r,\mathbf r')
=
-4\pi\delta(\mathbf r-\mathbf r').
\label{eq:Green_function_Helmholtz}
\end{equation}

La segunda identidad de Green aplicada a las funciones \(U(\mathbf r')\) y \(G(\mathbf r,\mathbf r')\) establece que

\begin{equation}
\int_V
\left[
G(\mathbf r,\mathbf r')\nabla'^2 U(\mathbf r')
-
U(\mathbf r')\nabla'^2 G(\mathbf r,\mathbf r')
\right]dV'
=
\oint_{\Sigma(V)}
\left[
G(\mathbf r,\mathbf r')
\frac{\partial U(\mathbf r')}{\partial n'}
-
U(\mathbf r')
\frac{\partial G(\mathbf r,\mathbf r')}{\partial n'}
\right]
d\sigma' .
\label{eq:second_Green_identity}
\end{equation}

En el integrando del lado izquierdo se puede sumar y restar el término \(k^2G(\mathbf r,\mathbf r')U(\mathbf r')\). Como \(k^2G U-k^2U G=0\), se obtiene

\begin{equation}
G\nabla'^2U-U\nabla'^2G
=
G(\nabla'^2+k^2)U
-
U(\nabla'^2+k^2)G.
\label{eq:Green_identity_Helmholtz_operator}
\end{equation}

Por tanto, la segunda identidad de Green puede escribirse directamente en términos del operador de Helmholtz como

\begin{equation}
\int_V
\left[
G(\mathbf r,\mathbf r')
(\nabla'^2+k^2)U(\mathbf r')
-
U(\mathbf r')
(\nabla'^2+k^2)G(\mathbf r,\mathbf r')
\right]dV'
=
\oint_{\Sigma(V)}
\left[
G\frac{\partial U}{\partial n'}
-
U\frac{\partial G}{\partial n'}
\right]
d\sigma' .
\label{eq:second_Green_identity_Helmholtz_operator}
\end{equation}

Ahora bien, como \(U\) satisface la ecuación homogénea de Helmholtz,

\begin{equation}
(\nabla'^2+k^2)U(\mathbf r')=0,
\end{equation}

mientras que \(G\) satisface

\begin{equation}
(\nabla'^2+k^2)G(\mathbf r,\mathbf r')
=
-4\pi\delta(\mathbf r-\mathbf r'),
\end{equation}

el miembro volumétrico se reduce a

\begin{equation}
\int_V
4\pi U(\mathbf r')\delta(\mathbf r-\mathbf r')\,dV'.
\end{equation}

Si el punto de observación \(\mathbf r\) pertenece al volumen \(V\), la propiedad de selección de la delta de Dirac da

\begin{equation}
\int_V
4\pi U(\mathbf r')\delta(\mathbf r-\mathbf r')\,dV'
=
4\pi U(\mathbf r).
\end{equation}

Se obtiene entonces

\begin{equation}
4\pi U(\mathbf r)
=
\oint_{\Sigma(V)}
\left[
G(\mathbf r,\mathbf r')
\frac{\partial U(\mathbf r')}{\partial n'}
-
U(\mathbf r')
\frac{\partial G(\mathbf r,\mathbf r')}{\partial n'}
\right]
d\sigma',
\end{equation}

y, finalmente,

\begin{equation}
U(\mathbf r)
=
\frac{1}{4\pi}
\oint_{\Sigma(V)}
\left[
G(\mathbf r,\mathbf r')
\frac{\partial U(\mathbf r')}{\partial n'}
-
U(\mathbf r')
\frac{\partial G(\mathbf r,\mathbf r')}{\partial n'}
\right]
d\sigma' .
\label{eq:Green_integral_theorem}
\end{equation}

Esta expresión es el teorema integral de Green para la ecuación de Helmholtz. En ella,

\begin{equation}
\frac{\partial}{\partial n'}
=
\mathbf n'\cdot\nabla'
\end{equation}

denota la derivada normal exterior sobre la frontera \(\Sigma(V)\).

Para aplicar esta representación al problema de propagación, se escoge una superficie cerrada formada por dos partes:

\begin{equation}
\Sigma(V)=\Sigma_1\cup\Sigma_2,
\end{equation}

donde \(\Sigma_1\) es el plano de apertura y \(\Sigma_2\) es una superficie hemisférica que cierra el volumen y cuyo radio se hace tender a infinito. De esta forma, la integral sobre la frontera se separa como

\begin{equation}
\begin{aligned}
U(\mathbf r)
=
\frac{1}{4\pi}
\int_{\Sigma_1}
\left[
G\frac{\partial U}{\partial n'}
-
U\frac{\partial G}{\partial n'}
\right]
d\sigma'
+
\frac{1}{4\pi}
\int_{\Sigma_2}
\left[
G\frac{\partial U}{\partial n'}
-
U\frac{\partial G}{\partial n'}
\right]
d\sigma' .
\end{aligned}
\label{eq:Green_integral_split_boundary}
\end{equation}

Considerando que el campo y la función de Green satisfacen la condición de radiación de Sommerfeld \cite{jackson1999}, la contribución sobre la superficie hemisférica \(\Sigma_2\) se anula en el límite en que su radio tiende a infinito:

\begin{equation}
\lim_{R\to\infty}
\int_{\Sigma_2}
\left[
G\frac{\partial U}{\partial n'}
-
U\frac{\partial G}{\partial n'}
\right]
d\sigma'
=
0.
\label{eq:Sommerfeld_condition_boundary_vanishes}
\end{equation}

Por consiguiente, la representación integral se reduce únicamente a la contribución sobre el plano de apertura \(\Sigma_1\):

\begin{equation}
U(\mathbf r)
=
\frac{1}{4\pi}
\int_{\Sigma_1}
\left[
G(\mathbf r,\mathbf r')
\frac{\partial U(\mathbf r')}{\partial n'}
-
U(\mathbf r')
\frac{\partial G(\mathbf r,\mathbf r')}{\partial n'}
\right]
d\sigma' .
\label{eq:Green_integral_aperture_plane}
\end{equation}

Esta es la forma espacial de la propagación escalar en términos del teorema integral de Green. El campo en el punto de observación \(\mathbf r\) queda determinado por los valores del campo \(U\) y de su derivada normal sobre el plano de apertura. Así, el problema de propagación se reformula como un problema de frontera: el campo no se obtiene mediante una descomposición espectral, sino mediante la contribución integral de los datos definidos sobre la superficie desde la cual emerge la onda.




Al elegir una función de Green $G(\mathbf{r}, \mathbf{r}')$ construida mediante el método de imágenes para que cumpla condiciones de frontera adecuadas (haciendo que la propia función de Green se anule en el plano), obtenemos:

\begin{equation}
G_1 = \frac{e^{ik\|\mathbf{r}-\mathbf{r}^*\|}}{\|\mathbf{r}-\mathbf{r}^*\|} - \frac{e^{ik\|\mathbf{r}-\mathbf{r}'\|}}{\|\mathbf{r}-\mathbf{r}'\|}.
\end{equation}

Reemplazando esta función y calculando su derivada normal respecto a la superficie, la expresión se reduce a la primera fórmula de difracción de Rayleigh-Sommerfeld:

\begin{equation}
U(\mathbf{r}) = -\frac{1}{2\pi} \int_{\Sigma} U(\mathbf{r}') \cos(\mathbf{n}, \mathbf{r}-\mathbf{r}') \left\{ ik - \frac{1}{\|\mathbf{r}-\mathbf{r}'\|} \right\} \frac{e^{ik\|\mathbf{r}-\mathbf{r}'\|}}{\|\mathbf{r}-\mathbf{r}'\|} d\sigma'.
\end{equation}

Dado que $\cos(\mathbf{n}, \mathbf{r}-\mathbf{r}') = \frac{z-z'}{\|\mathbf{r}-\mathbf{r}'\|}$, reescribimos:

\begin{equation}
U(\mathbf{r}) = -\frac{1}{2\pi} \int_{\Sigma} U(\mathbf{r}') \frac{z-z'}{\|\mathbf{r}-\mathbf{r}'\|} \left\{ ik - \frac{1}{\|\mathbf{r}-\mathbf{r}'\|} \right\} \frac{e^{ik\|\mathbf{r}-\mathbf{r}'\|}}{\|\mathbf{r}-\mathbf{r}'\|^2} d\sigma'.
\end{equation}

Hasta aquí, los desarrollos han sido exactos. La primera aproximación física que vamos a aplicar consiste en analizar el término entre llaves:

\begin{equation}
ik - \frac{1}{\|\mathbf{r}-\mathbf{r}'\|} = \frac{2\pi i}{\lambda} - \frac{1}{\|\mathbf{r}-\mathbf{r}'\|}.
\end{equation}

Bajo la justificación de que estudiaremos la propagación a distancias $\|\mathbf{r}-\mathbf{r}'\| \gg \lambda$, el primer término representa una contribución mucho mayor y domina completamente. Despreciando el segundo término, la integral toma la forma:

\begin{equation}
U(\mathbf{r}) = -\frac{ik}{2\pi} (z-z') \int_{\Sigma} U(\mathbf{r}') \frac{e^{ik\|\mathbf{r}-\mathbf{r}'\|}}{\|\mathbf{r}-\mathbf{r}'\|^2} d\sigma',
\end{equation}
que, reordenando los factores de número de onda, resulta en la expresión clásica del principio de Huygens:

\begin{equation}
U(\mathbf{r}) = \frac{(z-z')}{i\lambda} \int_{\Sigma} U(\mathbf{r}') \frac{e^{ik\|\mathbf{r}-\mathbf{r}'\|}}{\|\mathbf{r}-\mathbf{r}'\|^2} d\sigma'.
\label{eq:huygens_principle}
\end{equation}

La asunción de que $\|\mathbf{r}-\mathbf{r}'\| \gg \lambda$ es sumamente plausible en la inmensa mayoría de escenarios ópticos macroscópicos, por lo que la ecuación \eqref{eq:huygens_principle} posee un extenso rango de validez.


\subsection{Aproximación de Fresnel}

Hasta ahora los métodos de propagación presentados han sido ``exactos'' al menos en el sentido de que resuelven exactamente la ecuación de Helmholtz. 

En la práctica resulta muy conveniente considerar que las distancias transversales al eje óptico (tanto en la apertura como en el plano de observación) son pequeñas comparadas con la distancia de propagación $z$.

De tal manera que en \eqref{eq:huygens_principle} podamos aproximar la distancia entre un punto de la imagen y un punto del objeto, que aparece en la amplitud según


\begin{equation}
\|\mathbf{r}-\mathbf{r}'\| \approx z-z'.
\end{equation}






El otro término $\|\mathbf{r}-\mathbf{r}'\|$ está presente en la fase de la onda, y la fase es más sensible que la amplitud, al estar multiplicada por el número de onda $k = 2\pi/\lambda$ que relativiza la distancia entre los puntos objeto e imagen respecto a la longitud de onda, es el número de veces que la longitud de onda debe cubrir para llegar al valor numérico de la distancia mencionada. 

Con esto en mente, refinamos nuestra aproximación de la distancia entre los puntos objeto e imagen hasta orden $4$ 



$$\|\mathbf{r}-\mathbf{r}'\| = d \left[ 1 + \frac{(x-x')^2}{2d^2} + \frac{(y-y')^2}{2d^2} + \mathscr O\left( \left( \frac{\| \vb r - \vb r' \|}{d}\right)^4 \right) \right]$$



$$\|\mathbf{r}-\mathbf{r}'\| \approx d \left[ 1 + \frac{(x-x')^2}{2d^2} + \frac{(y-y')^2}{2d^2} \right]$$


donde $d=z-z'$.



Sustituyendo estas aproximaciones en la integral de propagación, obtenemos la conocida \textbf{Integral de Fresnel}:


\subsection{Integral de Fresnel}



\begin{equation}
U(x,y) = \frac{e^{ikz}}{i\lambda z} \iint_{-\infty}^{\infty} U(x',y') e^{i\frac{k}{2z} \left[ (x-x')^2 + (y-y')^2 \right]} dx'dy'.
\label{eq:fresnel_integral}
\end{equation}

Matemáticamente, este resultado revela una naturaleza lineal e invariante ante desplazamientos espaciales, lo cual permite escribir la propagación de manera directa en términos de una convolución espacial:

\begin{equation}
U(x,y) = \frac{e^{ikz}}{i\lambda z} \left( U(x,y) * e^{i\frac{k}{2z}(x^2+y^2)} \right).
\end{equation}

Asimismo, si expandimos el término cuadrático dentro de la exponencial en la integral de Fresnel \eqref{eq:fresnel_integral}:
\begin{equation*}
(x-x')^2 + (y-y')^2 = (x^2 + y^2) + (x'^2 + y'^2) - 2(xx' + yy'),
\end{equation*}
podemos reescribir el campo de forma equivalente extrayendo fuera de la integral los factores que no dependen de $(x', y')$:

\begin{equation}
U(x,y;z) = \frac{e^{ikz}}{i\lambda z} e^{i\frac{k}{2z}(x^2+y^2)} \iint_{-\infty}^{\infty} \left[ U(x',y') e^{i\frac{k}{2z}(x'^2+y'^2)} \right] e^{-i\frac{2\pi}{\lambda z}(xx'+yy')} dx'dy'.
\end{equation}

La integral resultante no es más que la Transformada de Fourier del campo de entrada modificado por una fase cuadrática. De esta manera, el campo en $z$ queda expresado mediante la operación:

\begin{equation}
U(x,y;z) = \frac{e^{ikz}}{i\lambda z} e^{i\frac{k}{2z}(x^2+y^2)} \mathscr{F} \left\{ e^{i\frac{k}{2z}(x'^2+y'^2)} U(x',y') \right\}(f_x, f_y),
\end{equation}
donde las frecuencias espaciales se evalúan específicamente en:
\begin{equation}
f_x = \frac{x}{\lambda z}, \quad f_y = \frac{y}{\lambda z}.
\end{equation}

Esta demostración conecta analíticamente la integral espacial de difracción con el uso fundamental de herramientas en el dominio de las frecuencias, cerrando así la conceptualización paralela entre este método paraxial y el método de espectro angular (ASM).



\section{Transformación que sufre el campo eléctrico al cambiar de medio} 


Considérese dos medios dieléctricos $M_1$ \textit{\&} $M_2$  distintos separados por una superficie $\Sigma$ cuya normal en cada punto se considera positiva cuando apunta hacia el medio $M_1$. 

Sean $n_1$ y $n_2$ los índices de refracción de $M_1$ y $M_2$ respectivamente, el campo eléctrico $\vb E_1$ en el medio $M_1$ y el campo eléctrico en el medio $M_2$ satisfacen


\begin{equation}
  \nabla^2 \vb E_i + k_i^2 \vb E_i = \vb 0,
  \qquad
  i=1,2.
\end{equation}


con 

\[
  k_i = n_i k_0 = \frac{2\pi n_i}{\lambda_0},
\]



Como la ecuación de Helmholtz a satisfacer es distinta en cada medio, esto genera una discontinuidad que transforma el campo que incide en la interfaz que separa ambos medios.


Escribiendo el campo por su espectro 

\begin{equation*} 
  \vb{E}(\vb x) = \frac{1}{(2\pi)^3} \int_{\mathbb K^3} \widehat{\vb{E}} (\vb{k}) e^{i\vb k \cdot \vb x} d^3k,
\end{equation*}


que por cumplir la ecuación de Helmholtz, es equivalente a 



\begin{equation*} 
  \vb{E}(\vb x) = \frac{1}{(2\pi)^3} \int_{\tilde{S}^2} \widehat{\vb{E}} (\vb{k}) e^{ik \uvec k \cdot \vb x} d\tilde{\Omega},
\end{equation*}

donde $\tilde S^2:=\{\hat{\mathbf k}\in\mathbb R^3:\|\hat{\mathbf k}\|=1\}$ es la esfera unitaria en el espacio $\mathbf{k}$, y $d\tilde\Omega$ su elemento de superficie (ángulo sólido), con $d\tilde\Omega=\sin\tilde\theta\,d\tilde\theta\,d\tilde\varphi$ en coordenadas esféricas.

Al remplazarlo en las condiciones de frontera, se encuentra la ley de transformación para cada componente $\vb k$ del campo. 



Dicha transformación puede ser efectuada de forma matemática mediante los coeficientes de Fresnel para una onda plana y monocromática, o la componente espectral de un campo general, la deducción de dicha transformación puede encontrarse en el apéndice \ref{sec:boundary_conditions_in_dielectrics}.


El campo incidente $\widehat{\vb E}(\vb k)$ se transforma de forma diferente en cada punto de la interfaz, y de manera diferente en dos polarizaciones principales $\uvec s$ y $\uvec p$.

Para definir las polarizaciones $\uvec s$ y $\uvec p$ primero hay que definir el plano $\Pi$ de incidencia

\begin{equation}
\Pi(\uvec{k},\uvec N) = \operatorname{span}\{\uvec{k}, \uvec{N}\},
\label{eq:pi}
\end{equation}

definimos así entonces los vectores unitarios $\uvec s$ y $\uvec p$ según

\begin{equation}
  \uvec{s}(\uvec{k},\uvec N)
  = \frac{\uvec N\times\uvec{k}}
  {\left\|\uvec N\times\uvec{k}\right\|},
\label{eq:sdef}
\end{equation}

\begin{equation}
  \uvec{p}(\uvec{k},\uvec N)
  = \uvec{s}(\uvec{k},\uvec N)\times\uvec{k}
  = \frac{(\uvec N\cdot\uvec{k})\uvec{k}-\uvec N}
  {\sqrt{1-\left(\uvec N\cdot\uvec{k}\right)^2}}.
\label{eq:pdef}
\end{equation}

El vector $\uvec{s}$ es perpendicular al plano de incidencia, mientras que $\uvec{p}$ pertenece a $\Pi(\uvec{k},\uvec N)$ y es perpendicular a $\uvec{k}$. Para el campo transmitido se definen análogamente $\uvec{s}'=\uvec{s}(\uvec{k}',\uvec N)$ y $\uvec{p}'=\uvec{p}(\uvec{k}',\uvec N)$.

\begin{figure}[H]
  \centering
  \thesisincludegraphics[width=0.65\textwidth]
    {figures/geometrical_setups/s_and_p_polarizations.pdf}
    {figures/fast/s_and_p_polarizations.jpg}
  \caption{Base de Fresnel en un punto $Q$ de la interfaz refractante. El plano de incidencia $\Pi = \operatorname{span}\{\hat{u}, \hat{N}\}$ contiene la normal $\hat{N}$ y la dirección de propagación incidente $\hat{u}$; el vector $\hat{p}$ (y su homólogo transmitido $\hat{p}'$) yace en $\Pi$ perpendicular a $\hat{u}$ (resp.\ $\hat{u}'$), mientras que $\hat{s}$ es perpendicular a $\Pi$.}
  \label{fig:fresnel_basis}
\end{figure}

Los coeficientes $t_s$, $t_p$ se conocen como coeficientes de Fresnel para la transmisión y se definen como

\begin{equation}
  t_s =
  \frac{2 n_1 \cos\theta_i}
  {n_1\cos\theta_i+n_2\cos\theta_t},
  \label{eq:ts}
\end{equation}

\begin{equation}
  t_p =
  \frac{2 n_1 \cos\theta_i}
  {n_2\cos\theta_i+n_1\cos\theta_t}.
  \label{eq:tp}
\end{equation}

Para escribirlos directamente en términos de la dirección del vector de onda incidente $\uvec{k}$, usamos

\begin{equation}
  \cos\theta_i = -\,\uvec N\cdot\uvec{k}.
  \label{eq:cos_i_kN}
\end{equation}

Por la ley de Snell, $n_1\sin\theta_i=n_2\sin\theta_t$, se tiene

\begin{equation}
  \cos\theta_t
  = \sqrt{1-\left(\frac{n_1}{n_2}\right)^2
  \left[1-\left(\uvec N\cdot\uvec{k}\right)^2\right]},
  \label{eq:cos_t_kN}
\end{equation}

donde se toma la rama positiva para la onda transmitida propagándose hacia el segundo medio. Sustituyendo \eqref{eq:cos_i_kN} y \eqref{eq:cos_t_kN} en \eqref{eq:ts} y \eqref{eq:tp}, los coeficientes de Fresnel de transmisión quedan como

\begin{equation}
t_s(\uvec{k},\uvec N) =
\frac{-2 n_1\,\uvec N\cdot\uvec{k}}
{-n_1\,\uvec N\cdot\uvec{k}
+ n_2\sqrt{1-\left(\frac{n_1}{n_2}\right)^2
\left[1-\left(\uvec N\cdot\uvec{k}\right)^2\right]}}
\label{eq:ts_kN}
\end{equation}

\begin{equation}
t_p(\uvec{k},\uvec N) =
\frac{-2 n_1\,\uvec N\cdot\uvec{k}}
{-n_2\,\uvec N\cdot\uvec{k}
+ n_1\sqrt{1-\left(\frac{n_1}{n_2}\right)^2
\left[1-\left(\uvec N\cdot\uvec{k}\right)^2\right]}}.
\label{eq:tp_kN}
\end{equation}



Con esto, conociendo el campo incidente sobre la interfaz $\Sigma$ que separa dos medios, podemos determinar el campo transmitido punto a punto sobre la misma interfaz. Si $\vb x_\Sigma\in\Sigma$, entonces la forma operacional de esta transformación es

\begin{equation}
  \vb{E}'_{\Sigma}(\vb x_\Sigma)
  =
  \mathscr{T}\!\left[
  \uvec{k},\uvec N(\vb x_\Sigma)
  \right]
  \left\{\vb{E}_{\Sigma}(\vb x_\Sigma)\right\}.
\end{equation}

En notación bra-ket, el operador local de transmisión se escribe como

\begin{equation}
  \mathscr{T}\!\left[
  \uvec{k},\uvec N(\vb x_\Sigma)
  \right]
  =
  t_p\!\left[
  \uvec{k},\uvec N(\vb x_\Sigma)
  \right]
  \ket{\uvec{p}'}
  \bra{\uvec{p}}
  +
  t_s\!\left[
  \uvec{k},\uvec N(\vb x_\Sigma)
  \right]
  \ket{\uvec{s}'}
  \bra{\uvec{s}}.
\end{equation}

De forma explícita, el campo transmitido queda dado por

\begin{equation}
  \begin{aligned}
  \vb{E}'_{\Sigma}(\vb x_\Sigma)
  &=
  t_p\!\left[
  \uvec{k},\uvec N(\vb x_\Sigma)
  \right]
  \left[
  \vb{E}_{\Sigma}(\vb x_\Sigma)\cdot
  \uvec{p}\!\left[
  \uvec{k},\uvec N(\vb x_\Sigma)
  \right]
  \right]
  \uvec{p}'\!\left[
  \uvec{k}',\uvec N(\vb x_\Sigma)
  \right]
  \\[4pt]
  &\quad+
  t_s\!\left[
  \uvec{k},\uvec N(\vb x_\Sigma)
  \right]
  \left[
  \vb{E}_{\Sigma}(\vb x_\Sigma)\cdot
  \uvec{s}\!\left[
  \uvec{k},\uvec N(\vb x_\Sigma)
  \right]
  \right]
  \uvec{s}'\!\left[
  \uvec{k}',\uvec N(\vb x_\Sigma)
  \right].
  \end{aligned}
  \label{eq:Transmited_E_field_in_Sigma}
\end{equation}

\textit{(véase \ref{sec:appendix_maxwell_frontera_fresnel} para una deducción detallada de este resultado)}.










En el problema \textit{estigmático simple} la dirección de propagación en función de cualquier punto $Q \ \in \Sigma$ es siempre conocida, pues por la construcción del problema estigmático


$$
\uvec k := \uvec u = \text{dirección}(A, Q),
$$


$$
\uvec k':= \uvec u' = \text{dirección}(Q, A'),
$$

donde dirección($AQ$) es el vector unitario que va desde $A$ hasta $Q$. 

Esto porque todos los rayos que emergen de $A$ lo hacen radialmente, e igualmente radialmente convergen en $A'$. 


Tenemos pues, todo lo necesario para pasar al estudio del campo eléctrico en el plano focal en un sistema simple refractivo.



% chapter  (end)
